\documentclass[11pt,a4paper]{article}

% Template visual Matriz Aprovação. Compile com XeLaTeX ou LuaLaTeX.
\usepackage{fontspec}
\usepackage{polyglossia}
\setmainlanguage{portuguese}
\IfFontExistsTF{Aptos}{
  \setmainfont{Aptos}
  \setsansfont{Aptos}
}{
  \IfFontExistsTF{Calibri}{
    \setmainfont{Calibri}
    \setsansfont{Calibri}
  }{
    \IfFontExistsTF{TeX Gyre Heros}{
      \setmainfont{TeX Gyre Heros}
      \setsansfont{TeX Gyre Heros}
    }{
      \setmainfont{Latin Modern Sans}
      \setsansfont{Latin Modern Sans}
    }
  }
}
\usepackage[a4paper,margin=2cm,headheight=16pt]{geometry}
\usepackage{graphicx}
\usepackage{xcolor}
\usepackage{tikz}
\usepackage{tcolorbox}
\usepackage{enumitem}
\usepackage{booktabs}
\usepackage{tabularx}
\usepackage{amssymb}
\usepackage{fancyhdr}
\usepackage{titlesec}
\usepackage{hyperref}

\definecolor{matrizlime}{HTML}{CBFF4D}
\definecolor{matrizink}{HTML}{101313}
\definecolor{matrizmuted}{HTML}{596363}
\definecolor{matrizline}{HTML}{DDE3DF}
\definecolor{matrizsoft}{HTML}{F3F7EC}

\hypersetup{colorlinks=true,linkcolor=matrizink,urlcolor=matrizink}
\pagestyle{fancy}
\fancyhf{}
\lhead{\textsf{\small MATRIZ APROVAÇÃO}}
\rhead{\textsf{\small APOSTILA}}
\cfoot{\textsf{\small \thepage}}
\renewcommand{\headrulewidth}{0.4pt}
\renewcommand{\headrule}{\hbox to\headwidth{\color{matrizline}\leaders\hrule height \headrulewidth\hfill}}

\titleformat{\section}{\Large\bfseries\color{matrizink}}{\thesection}{0.6em}{}
\titleformat{\subsection}{\large\bfseries\color{matrizink}}{\thesubsection}{0.6em}{}
\setlist{nosep,leftmargin=*}

\newtcolorbox{foco}{colback=matrizsoft,colframe=matrizlime,boxrule=1pt,arc=3pt,left=10pt,right=10pt,top=8pt,bottom=8pt}
\newtcolorbox{lei}{colback=matrizink,colframe=matrizink,coltext=white,boxrule=0pt,arc=3pt,left=10pt,right=10pt,top=8pt,bottom=8pt}
\newtcolorbox{alerta}{colback=white,colframe=matrizline,boxrule=0.6pt,arc=3pt,left=10pt,right=10pt,top=8pt,bottom=8pt}

% Logo usada na capa. Caminho relativo à pasta onde o .tex é compilado
% (materiais/<disciplina>/). Ajuste se a estrutura mudar.
\newcommand{\logomatriz}{../../public/logo.png}

\begin{document}

\begin{titlepage}
  \thispagestyle{empty}
  \begin{center}
    \includegraphics[width=0.55\linewidth]{\logomatriz}\par
  \end{center}
  \vspace{2.2cm}
  {\sffamily\fontsize{28}{32}\selectfont\bfseries Apostila — Direito Administrativo: Princípios da Administração Pública\par}
  \vspace{0.4cm}
  {\sffamily\large Receita Federal, CGU, PRF, Polícia Federal, INSS, TRT, TCU/TCE, Câmara e OAB\par}
  \vfill
  \begin{foco}
    \textsf{\textbf{Foco da apostila}}\\[3pt]
    
  \end{foco}
  \vspace{0.7cm}
  {\sffamily\small Atualizado em 16 de agosto de 2026\quad ·\quad Cebraspe, FGV e FCC\par}
  \vspace{1cm}
  {\sffamily\small matrizaprova.com\par}
\end{titlepage}

\tableofcontents
\newpage

% Conteúdo gerado entra aqui.
\section*{O que cai}

O núcleo mais cobrado é o conjunto de princípios expressos no art. 37, caput: \textbf{legalidade, impessoalidade, moralidade, publicidade e eficiência}. A memorização tradicional é \textbf{LIMPE}.

Também aparecem com frequência:

\begin{enumerate}
  \item diferença entre legalidade para o particular e legalidade administrativa;
  \item impessoalidade, finalidade pública e vedação à promoção pessoal;
  \item moralidade como princípio jurídico, não apenas ético;
  \item publicidade como regra e sigilo como exceção constitucional;
  \item eficiência sem autorização para violar a lei;
  \item aplicação dos princípios à Administração direta e indireta de todos os
\end{enumerate}

   Poderes e entes federativos.

\section*{Teoria essencial}

\subsection*{1. Abrangência do art. 37}

O caput do art. 37 determina que a Administração Pública direta e indireta de qualquer dos Poderes da União, dos Estados, do Distrito Federal e dos Municípios obedecerá aos princípios expressos no dispositivo.

Isso significa que os princípios se aplicam:

\begin{itemize}
  \item ao Executivo, Legislativo e Judiciário, quando exercem função administrativa;
  \item à União, aos Estados, ao Distrito Federal e aos Municípios;
  \item à Administração direta e à indireta;
  \item às autarquias, fundações públicas, empresas públicas e sociedades de economia mista, observadas as particularidades de cada entidade.
\end{itemize}

Não é correto limitar o art. 37 ao Poder Executivo. Um tribunal, por exemplo, também deve respeitar os princípios quando realiza licitação, nomeia servidor ou administra seu orçamento.

\subsection*{2. Legalidade administrativa}

Para o particular, vale a ideia de que ninguém é obrigado a fazer ou deixar de fazer algo senão em virtude de lei. Na Administração Pública, a legalidade tem sentido mais rigoroso: o agente público só pode agir quando houver autorização ou fundamento jurídico que permita sua atuação.

Assim:

\begin{itemize}
  \item o particular pode fazer tudo que a lei não proíbe;
  \item a Administração só pode fazer o que a ordem jurídica autoriza.
\end{itemize}

Legalidade não significa obediência cega a qualquer ordem. O ato administrativo deve respeitar a Constituição, as leis, os regulamentos pertinentes e os demais princípios jurídicos.

\subsection*{3. Impessoalidade}

A impessoalidade possui aplicações relacionadas entre si:

\subsubsection*{3.1. Finalidade pública}

O agente deve atuar para atender ao interesse público previsto na norma, e não para favorecer ou prejudicar uma pessoa determinada. O ato praticado com finalidade diferente da prevista caracteriza desvio de finalidade.

\subsubsection*{3.2. Igualdade de tratamento}

A Administração deve tratar situações equivalentes de maneira equivalente, sem discriminações arbitrárias. Diferenças de tratamento podem existir quando há justificativa jurídica e relação com a finalidade da medida.

\subsubsection*{3.3. Imputação do ato ao órgão}

A atuação administrativa é atribuída ao órgão ou à entidade pública, não à figura pessoal do agente. O servidor manifesta a vontade administrativa dentro de sua competência.

\subsubsection*{3.4. Vedação à promoção pessoal}

O art. 37, § 1º, determina que a publicidade dos atos, programas, obras, serviços e campanhas dos órgãos públicos deve ter caráter educativo, informativo ou de orientação social. Não pode conter nomes, símbolos ou imagens que caracterizem promoção pessoal de autoridades ou servidores.

Uma campanha pública pode informar quem presta o serviço, mas não pode usar a comunicação institucional para construir propaganda pessoal do governante.

\subsection*{4. Moralidade}

Moralidade administrativa exige comportamento ético, honesto, leal e compatível com a finalidade pública. Não se resume à moral privada do agente.

É princípio jurídico e pode fundamentar o controle de atos administrativos. Mesmo quando um comportamento parece formalmente autorizado, pode ser inválido se contrariar padrões jurídicos de probidade, lealdade e boa-fé.

A moralidade não autoriza o intérprete a substituir a lei por preferências pessoais. Sua aplicação deve ser fundamentada no sistema jurídico, nos fatos e na finalidade pública.

\subsection*{5. Publicidade}

Publicidade é a divulgação dos atos administrativos para permitir transparência, controle social e produção de efeitos quando exigida pelo ordenamento.

Ela não significa que toda informação pública seja irrestrita. A Constituição admite restrições quando o sigilo for imprescindível à segurança da sociedade e do Estado, além da proteção da intimidade, vida privada, honra e imagem.

Publicidade também não equivale a propaganda pessoal. A divulgação deve ser institucional e cumprir finalidade educativa, informativa ou de orientação social.

\subsection*{6. Eficiência}

Eficiência exige atuação administrativa com qualidade, produtividade, economicidade e redução de desperdícios, buscando bons resultados para a coletividade.

O princípio não permite:

\begin{itemize}
  \item descumprir procedimento obrigatório;
  \item ignorar direitos do administrado;
  \item escolher uma solução ilegal por ser mais rápida;
  \item substituir critérios jurídicos por metas meramente numéricas.
\end{itemize}

A Administração deve buscar o melhor resultado \textbf{dentro da legalidade}. A eficiência complementa os demais princípios; não os elimina.

\subsection*{7. Princípios expressos e implícitos}

O LIMPE não é uma lista fechada. A Administração também se submete a princípios reconhecidos pela Constituição, pelas leis e pela interpretação do Direito, como supremacia do interesse público, indisponibilidade do interesse público, razoabilidade, proporcionalidade, motivação, segurança jurídica, continuidade do serviço público e autotutela.

Princípio \textbf{expresso} é aquele escrito de forma direta na Constituição ou na lei. Princípio \textbf{implícito} é extraído do sistema jurídico, mesmo sem estar apresentado com essa denominação em um único dispositivo.

\subsection*{8. Regime jurídico e ponderação}

Os princípios devem ser aplicados em conjunto. Uma decisão eficiente, mas secreta sem justificativa, viola publicidade. Uma decisão publicizada, mas usada para promover o governante, viola impessoalidade. Uma decisão legal em sentido formal, mas desonesta e desviada de finalidade, pode violar moralidade.

Em caso de tensão, a solução exige análise do caso concreto, da finalidade constitucional e da proporcionalidade. Não existe hierarquia automática entre os cinco princípios do LIMPE.

\section*{Como a banca cobra}

\begin{itemize}
  \item \textbf{LIMPE:} legalidade, impessoalidade, moralidade, publicidade e eficiência.
  \item \textbf{Particular x Administração:} a Administração não tem a mesma liberdade de atuação do particular.
  \item \textbf{Publicidade:} é regra, mas admite restrições justificadas.
  \item \textbf{Eficiência:} não supera a legalidade.
  \item \textbf{Impessoalidade:} abrange finalidade pública, igualdade e vedação à promoção pessoal.
  \item \textbf{Moralidade:} é princípio jurídico controlável, não simples recomendação moral.
  \item \textbf{Abrangência:} art. 37 alcança todos os Poderes e entes federativos quando exercem função administrativa.
\end{itemize}

\section*{Exemplos resolvidos}

\subsection*{Exemplo 1}

Um prefeito divulga obra pública com seu nome, fotografia e slogan pessoal. A conduta é compatível com o art. 37, § 1º?

\textbf{Não.} A publicidade institucional pode ter caráter educativo, informativo ou de orientação social, mas não pode conter elementos de promoção pessoal de autoridade ou servidor.

\subsection*{Exemplo 2}

Um órgão decide ignorar a exigência legal de procedimento porque encontrou uma forma mais rápida de atender os cidadãos. A decisão é válida com fundamento no princípio da eficiência?

\textbf{Não.} Eficiência exige bons resultados, mas deve ser buscada dentro da legalidade. O princípio não transforma uma obrigação legal em opção.

\section*{Quadro-resumo}

\begin{tabularx}{\linewidth}{llX}
\toprule
 \textbf{Princípio} & \textbf{Palavra-chave} & \textbf{Pergunta de prova} \\
\midrule
 Legalidade & autorização jurídica & a Administração podia agir assim? \\
 Impessoalidade & finalidade & houve favorecimento, perseguição ou promoção pessoal? \\
 Moralidade & probidade & o comportamento é honesto e leal à finalidade pública? \\
 Publicidade & transparência & o ato foi divulgado ou o sigilo foi justificado? \\
 Eficiência & resultado & a atuação entrega qualidade sem violar a lei? \\
\bottomrule
\end{tabularx}


\section*{Questões de fixação}

\subsection*{Questão 1 — (questão inédita)}

São princípios expressos da Administração Pública previstos no art. 37, caput, da Constituição Federal:

A) supremacia, autotutela, continuidade, publicidade e eficiência. \\ B) legalidade, impessoalidade, moralidade, publicidade e eficiência. \\ C) proporcionalidade, razoabilidade, legalidade, moralidade e continuidade. \\ D) indisponibilidade, motivação, publicidade, eficiência e hierarquia. \\ E) legalidade, finalidade, segurança jurídica, publicidade e eficiência.

\begin{alerta}
\textbf{Gabarito: B.} O art. 37, caput, apresenta o conjunto tradicionalmente memorizado por LIMPE. Os demais princípios podem ser reconhecidos no regime administrativo, mas não compõem essa lista literal do caput.
\end{alerta}

\subsection*{Questão 2 — (questão inédita)}

Sobre o princípio da legalidade administrativa, é correto afirmar que:

A) o agente público pode fazer tudo que não for expressamente proibido. \\ B) a Administração pode agir sem previsão jurídica quando busca eficiência. \\ C) a Administração deve atuar nos limites e autorizações da ordem jurídica. \\ D) a legalidade se aplica apenas aos atos do Poder Executivo. \\ E) a lei não limita atos discricionários.

\begin{alerta}
\textbf{Gabarito: C.} A Administração está vinculada à ordem jurídica. A discricionariedade também possui limites legais e principiológicos.
\end{alerta}

\subsection*{Questão 3 — (questão inédita)}

A publicação de campanha institucional com nome e fotografia da autoridade responsável caracteriza possível violação principalmente ao princípio da:

A) eficiência. B) continuidade. C) impessoalidade. D) autotutela. E) hierarquia.

\begin{alerta}
\textbf{Gabarito: C.} O art. 37, § 1º, veda elementos de promoção pessoal na publicidade institucional.
\end{alerta}

\subsection*{Questão 4 — (questão inédita)}

O princípio da publicidade:

A) impede qualquer hipótese de sigilo administrativo. \\ B) exige divulgação de toda informação sem considerar direitos individuais. \\ C) é regra, mas pode admitir restrições constitucionalmente justificadas. \\ D) autoriza promoção pessoal em campanhas educativas. \\ E) aplica-se apenas a atos do Executivo.

\begin{alerta}
\textbf{Gabarito: C.} A transparência é regra, mas a Constituição admite hipóteses de sigilo e proteção da intimidade, entre outras restrições jurídicas.
\end{alerta}

\subsection*{Questão 5 — (questão inédita)}

Uma medida administrativa reduz custos, mas descumpre uma obrigação legal. À luz dos princípios constitucionais, a medida:

A) é válida porque a eficiência sempre prevalece. \\ B) é válida se gerar economia ao erário. \\ C) é inválida, pois eficiência não autoriza violação à legalidade. \\ D) é obrigatória quando aprovada pelo superior hierárquico. \\ E) é válida se não houver prejuízo financeiro.

\begin{alerta}
\textbf{Gabarito: C.} A busca por eficiência deve ocorrer dentro da Constituição e da lei. Resultado econômico não legitima atuação ilegal.
\end{alerta}

\section*{Revisão final}

\begin{itemize}
  \item[$\square$] Memorizei o LIMPE.
  \item[$\square$] Sei que o art. 37 alcança todos os Poderes quando exercem função administrativa.
  \item[$\square$] Diferencio legalidade do particular e legalidade administrativa.
  \item[$\square$] Relaciono impessoalidade a finalidade, igualdade e promoção pessoal.
  \item[$\square$] Reconheço moralidade como princípio jurídico.
  \item[$\square$] Sei que publicidade é regra, mas admite sigilo fundamentado.
  \item[$\square$] Sei que eficiência não permite descumprir a lei.
  \item[$\square$] Diferencio princípios expressos de princípios implícitos.
\end{itemize}

\end{document}
